\section{Regla de tres simple}
\subsection{Teoría}
Es la operación que se utiliza para encontrar un término en una relación, esto quiere decir que si a y b son dos números entonces:
\begin{equation*}
    \frac{a}{b}= \text{constante}
\end{equation*}
\begin{ejemplo}
    Si 12 discos compactos cuestan \$600 ¿cuánto costarán 18?\\
    Las cantidades son directamente proporcionales, ya que al aumentar el número de discos el precio también se
    incrementa. Se forma una proporción entre las razones del supuesto y la pregunta.
    \begin{align*}
        \frac{12}{18} & = \frac{600}{x}       \\
        x             & =\frac{(600)(18)}{12} \\
        x             & = \frac{10800}{12}    \\
        x             & = 900
    \end{align*}
    Por tanto, 18 discos compactos cuestan \$900
\end{ejemplo}
lo importante de estos problemas es la organización de las cantidades que tenemos.
\begin{ejemplo}
    Una llave que se abre 4 horas diarias durante 5 días, vierte 5 200 litros de agua, ¿cuántos litros vertirá en 12 días si se abre 4 horas por día?\\
    Se calcula el número de horas totales; es decir, en 5 días la llave ha estado abierta 20 horas y en 12 días la llave permaneció abierta 48 horas.
    \begin{align*}
        \frac{20}{48} & = \frac{5200}{x}       \\
        x             & =\frac{(5200)(48)}{20} \\
        x             & = \frac{249600}{20}    \\
        x             & = 12480
    \end{align*}
    Por consiguiente, en 48 horas la llave vierte 12 480 litros.
\end{ejemplo}
\subsection{Ejercicios}
\begin{enumerate}
    \item Liam escucha la radio durante 30 minutos, lapso en el que hay 7 minutos de anuncios comerciales; si escucha la radio durante 120 minutos, ¿cuántos minutos de anuncios escuchará?
    \item Un leñador tarda 8 segundos en dividir en 4 partes un tronco de cierto tamaño ¿cuánto tiempo tardará en dividir un tronco semejante en 5 partes?
    \item El valor de 25 m 2 de azulejo es de \$3 125. ¿Cuántos m 2 se comprarán con \$15 625?
    \item Un microbús cobra a una persona \$17.50 de pasaje por una distancia de 21 kilómetros, ¿cuánto pagará otra persona, cuyo destino está a 51 kilómetros de distancia?
    \item Si 3 decenas de pares de zapatos cuestan \$18 000, ¿cuál será el precio de 25 pares?
    \item Don Arturo tiene una pastelería y sabe que para hacer un pastel de fresas para 8 personas utiliza 2 kg de azúcar, ¿qué cantidad de azúcar utilizará si le encargan un pastel, también de fresas, que alcance para 24 personas?
\end{enumerate}